\chapter{Uso dell'AI}\label{ch:ai}
Il capitolo distingue le dichiarazioni degli autori, le attività riscontrabili
nei materiali conservati e ciò che tali materiali non permettono di stabilire.
È completo e proposto agli autori per revisione, senza attribuirne
l'approvazione. La matrice affermazioni--fonti e i controlli documentali sono
in \code{docs/relazione/verifiche/ticket-07/}; l'integrazione del
riallineamento è documentata nelle verifiche del ticket~13.

\begingroup
\emergencystretch=6em
\def\UrlBreaks{\do\a\do\b\do\c\do\d\do\e\do\f\do\g\do\h\do\i\do\j\do\k\do\l\do\m\do\n\do\o\do\p\do\q\do\r\do\s\do\t\do\u\do\v\do\w\do\x\do\y\do\z\do\A\do\B\do\C\do\D\do\E\do\F\do\G\do\H\do\I\do\J\do\K\do\L\do\M\do\N\do\O\do\P\do\Q\do\R\do\S\do\T\do\U\do\V\do\W\do\X\do\Y\do\Z}

\section{Progettazione e sviluppo iniziale manuali}
% AI-01
Secondo la dichiarazione degli autori registrata nel piano, la progettazione,
la struttura iniziale del software e le prime classi sono state realizzate
manualmente; tutta la revisione è stata svolta nelle sessioni con Codex.
Questo è il significato attribuito allo \emph{scheletro} iniziale.
Si tratta di una dichiarazione di provenienza, non di un'attribuzione
ricavata autonomamente dal codice.\footnote{\raggedright \code{docs/relazione/piano.md},
sezione \emph{Origine del lavoro e uso dell'AI}. Tutti i percorsi citati
sono relativi alla radice del repository.}
La cronologia Git conserva versioni e modifiche, ma non separa per ogni
riga il lavoro manuale da quello assistito. Non si assegnano quindi singole
classi all'uno o all'altro sulla sola base di nomi, stile o autore del commit.

\section{Revisione assistita da Codex}
% AI-02
La fotografia dei registri conserva audit, piano e diciotto ticket della
prima fase. L'audit del 30 ottobre 2026 individua rischi su denaro, stati,
concorrenza e coerenza documentale; dichiara anche che PostgreSQL e display
non erano disponibili. I controlli con DAO simulati di quella fase non
dimostravano l'isolamento del database. Il piano trasforma i rilievi in
decisioni, dipendenze e criteri verificabili; i ticket e le evidenze
registrano implementazione, tentativi falliti e chiusure. L'attribuzione
complessiva a Codex resta quella dichiarata dagli autori.\footnote{\raggedright 
\code{docs/evidence/revision-records.zip}: \code{AUDIT_DRIVEHUB_2026-10-30.md},
\emph{Verifiche svolte}; \code{PIANO_REVISIONI_DRIVEHUB.md},
\emph{Decisioni raccolte} e \emph{Attività e dipendenze};
\code{TRACKER_DRIVEHUB/revisioni/issues/}. Il manifest interno
\code{SHA256.json} identifica la fotografia, non una trascrizione delle sessioni.}

% AI-03
Un esempio è il prezzo manuale. La proposta dell'audit di eliminarlo viene
superata dalla decisione degli autori: mantenere una base scelta dal manager
e applicarle la formula comune. La revisione raggiunge dominio, servizi,
DAO e viste; una regressione documenta base 10.000 e sconto 10\% con risultato
9.000. La review Spec rileva poi una differenza sui centesimi prima e dopo
la persistenza, riprodotta e corretta con normalizzazione condivisa.
Un intervento distinto fissa il totale all'acconto: 20.000, acconto 4.000,
saldo 16.000 anche dopo cambi di listino. Sono esempi di requisito deciso,
proposta tecnica, controprova e correzione, non prove dell'origine di ogni
riga.\footnote{\raggedright \code{docs/adr/0001-base-manuale-e-prezzo-concordato.md};
\code{docs/evidence/r03-prices/verification.md}, \emph{TDD e verifiche}
e \emph{Revisione indipendente};
\code{docs/evidence/r04-agreement/verification.md}, \emph{TDD e comandi}.}

% AI-04
Per le tabelle JavaFX, le tracce mostrano letture collettive e caricamenti
asincroni, con controlli di rendering senza query, letture bloccate e risultati
obsoleti. La pulizia successiva rimuove soltanto componenti senza chiamanti.
La revisione comprende inoltre identità, rollback e gare sui flussi economici.
Interventi successivi alla prima fase riallineano richieste e registro al
punto di vista della concessionaria e uniformano gli enum PostgreSQL:
i vecchi nomi e segni nelle prove rimangono riferiti alla loro versione,
come precisato nell'appendice~\ref{app:verifiche}.\footnote{\raggedright 
\code{docs/evidence/r08-tables/verification.md};
\code{docs/evidence/r10-cleanup/verification.md};
\code{docs/evidence/r11-integrated/acceptance-matrix.md};
\code{docs/evidence/request-semantics-01/verification.md};
\code{docs/evidence/sale-semantics-02/verification.md};
\code{docs/evidence/postgres-enums-2026-11-09/verification.md}.}

\section{Verifica e controllo umano}
% AI-05
Le decisioni umane documentate riguardano requisiti e perimetro: l'ADR~0001
registra la scelta della base manuale e del totale concordato; l'ADR~0002
registra la scelta degli enum nativi e le condizioni dell'aggiornamento
locale, con backup e prova su copia. Sono decisioni registrate come
confermate dall'utente; non dimostrano una lettura umana di ciascuna patch
o l'approvazione di ogni risultato.\footnote{\raggedright ADR~0001 citato sopra;
\code{docs/adr/0002-enumerazioni-native-postgresql.md}.
Per gli esiti storici di conservazione e ripristino:
\code{docs/evidence/db-consolidation-2026-11-09/verification.md}
e \code{docs/evidence/postgres-enums-2026-11-09/verification.md}.
Nessun backup privato è stato consultato per questo capitolo.}

% AI-06
Le verifiche tecniche sono un livello distinto. Log, XML, manifest e patch
permettono di controllare risultati e versione esaminata. Ad esempio R11
registra il 3 novembre 394 prove riuscite su JDK~21 e PostgreSQL~16 dedicato;
il 9 novembre il log degli enum ne registra 423. Non si sommano sessioni
diverse né si ricava da questi conteggi una percentuale di copertura.
Il secondo risultato non dispone di manifest runtime o XML per metodo:
il rapporto con il codice corrente è ricostruito tramite fonti e Git.
Il capitolo~\ref{ch:testing} e l'appendice~\ref{app:verifiche} distinguono
questi limiti e le prove JavaFX. Qui non sono stati rieseguiti test
applicativi, harness o accessi al database.

% AI-07
Anche le review Standards e Spec conservate sono revisioni assistite,
svolte da agenti in contesti separati: i loro esiti non sono approvazioni
umane. Per la relazione, le verifiche 05--06 documentano correzioni,
compilazione e ispezione del PDF; i capitoli sono stati proposti agli
autori e restano da loro rivedere. L'eventuale approvazione non viene
dedotta dal silenzio o dalla chiusura di un ticket. Anche le review separate
dei ticket documentali 09--12 chiudono valutazioni tecniche, comprese le
correzioni ricontrollate, senza attestare una validazione degli autori.\footnote{\raggedright 
\code{docs/evidence/r03-prices/verification.md}, \emph{Revisione indipendente};
\code{docs/relazione/verifiche/ticket-05/revisione.md} e
\code{docs/relazione/verifiche/ticket-06/revisione.md};
\code{docs/relazione/verifiche/ticket-09/revisione.md},
\code{ticket-10/revisione.md}, \code{ticket-11/revisione.md} e
\code{ticket-12/revisione.md}, sotto la stessa directory delle verifiche.}

\section{Documentazione, diagrammi e limiti}
% AI-08
Il supporto di Codex comprende la documentazione: il README lo cita
esplicitamente per l'analisi dell'avvio, lo script Windows e le istruzioni.
Le evidenze R09 registrano la revisione tecnica e teorica, distinguendo
factory statiche da Factory Method e meccanismi DAO/DI dai pattern GoF.
Registrano inoltre diagrammi PlantUML di quella fase, separati dagli
artefatti storici.\footnote{\raggedright \code{README.md}, \emph{Uso dell'AI};
\code{docs/evidence/r09-instructions/verification.md},
\emph{Cambiamenti e motivazione};
\code{docs/evidence/r09-diagrams/verification.md},
\emph{Diagrammi correnti e storico}.}

% AI-09
Anche questa riscrittura della relazione è assistita da Codex: comprende
sintesi dei sorgenti e delle evidenze, selezione dei frammenti, matrici di
tracciabilità, composizione LaTeX e controllo del PDF. Le nuove tavole della
relazione usano draw.io: sono viste selettive riscontrate su Java e schema,
non un'esportazione UML automatica del codice. Sorgenti modificabili,
esportazioni e verbali restano conservati. Questo capitolo riusa le tavole
del capitolo~2 e dell'appendice~\ref{app:classi}, senza crearne altre.
Le verifiche documentali 01--07 distinguono i controlli attuali dalle
prove applicative pregresse.\footnote{\raggedright \code{docs/relazione/diagrammi/README.md};
\code{docs/relazione/verifiche/ticket-03/matrice-copertura.md};
\code{docs/relazione/verifiche/ticket-04/README.md};
\code{docs/relazione/verifiche/ticket-05/README.md},
\code{docs/relazione/verifiche/ticket-06/README.md}
e \code{docs/relazione/verifiche/ticket-07/README.md}.}

Le figure~\ref{fig:uc-autori}, \ref{fig:nav-autori} e~\ref{fig:er-autori}
derivano invece dalle tre tavole realizzate dagli autori e da loro indicate
per l'integrazione. Codex ne ha corretto i contenuti superati e curato
l'inserimento, mantenendo stile e struttura; questa revisione non attribuisce
all'AI la progettazione delle tavole originali. Le copie corrette e i
controlli sono conservati in \code{docs/relazione/verifiche/tavole-autori/}.

% AI-11
Su richiesta degli autori registrata nel piano il 16 novembre 2026,
il riallineamento 09--13 conserva anche la scansione interna e la funzione
espositiva dei materiali di \code{Relazione_SWE}. Codex ha confrontato
le 102 voci e i 23 frammenti originali con spiegazioni, esempi e oracoli
correnti; le sei schede principali sono complete nel corpo e le aggiuntive
in A. Il ticket~10 documenta gli undici rifacimenti/adattamenti draw.io:
quattro ripartono dai nativi degli autori, sette dai riferimenti visivi
e PlantUML originali. Non si attribuiscono all'AI i disegni di partenza.
Il ticket~13 riusa questi artefatti, integra il documento e verifica
compilazione, tutte le pagine e autonomia dei sorgenti tramite estrazione
e ricompilazione dello ZIP. Baseline, corrispondenze, log, verbale visivo
e review separate identificano gli artefatti effettivamente controllati.
Sono verifiche documentali: non nuove esecuzioni Java, test, harness o DB,
né un'approvazione umana.\footnote{\raggedright
\code{docs/relazione/piano.md}, \emph{Riallineamento richiesto dagli autori};
\code{docs/relazione/allineamento.md};
\code{docs/relazione/verifiche/ticket-10/tavole.md};
\code{docs/relazione/verifiche/ticket-13/README.md} e
\code{docs/relazione/verifiche/ticket-13/capitolo-5-fonti.md}.}

% AI-10
Le tracce non sono una cronologia completa delle conversazioni. Non
consentono di ricostruire prompt integrali, versioni dei modelli, tempi
complessivi, percentuali di codice generato o una precisa modalità
d'integrazione nell'IDE. I metadati dei commit e le impronte dei file
identificano artefatti, non il loro autore intellettuale. I risultati
storici restano tali; compilazione, ispezione del PDF e review documentali
non certificano il software corrente e non sostituiscono il controllo
degli autori.
\endgroup
